% §15.5 Bridge to Part V
%
% Monograph register: close Part IV (pair BHF and BPHF as the two
% operating points on the perfect-filter line, with the chapter-8 floor),
% note the (eps,0)/(0,omega) endpoints with hybrids interpolating
% (salvaged ex-11), preview Part V (ch16 NOT-Bernoulli, ch17-18 cipher
% maps / encrypted search). Cuts the "what 15.6 supplies" subsection.
% Cross-refs def:one-sided, def:perfect-filter-line, thm:entropy-floor,
% def:bhf, thm:bhf-optimality, def:phf, prop:phf-space, def:bphf, prop:bphf-space, sec:ranked-search,
% sec:bib-notes-ch15. DEFINED: sec:bridge-to-part5.

\section{Bridge to Part V}
\label{sec:bridge-to-part5}

The chapter assembled the constructive answer to the FPR-zero arm of the
perfect-filter line. The BPHF of \Cref{def:bphf} reaches $\fprate = 0$ in
the witness query model, is exact at full resolution, and trades a
controlled $\fnrate \ge 1 - m/n$ for a smaller table through a
$k$-perfect hash, at $\lceil \log_2 n \rceil$ bits per element plus the
addressing cost of a minimal perfect hash (\Cref{prop:bphf-space}),
applying cleanly to ranked retrieval through the construction of
\Cref{sec:ranked-search}.

\subsection*{Closing Part IV}

Part~IV has a clean shape in retrospect. Chapters~5 and~6 named the
perfect-filter line (\Cref{def:perfect-filter-line}, building on\index{perfect filter line}
\Cref{def:one-sided}); chapter~8 derived the entropy floor
(\Cref{thm:entropy-floor}) as the bit-budget lower bound any FNR-zero
structure must respect; chapters~14 and~15 constructed the operating
points at the two extreme error axes. On the FNR-zero arm the BHF meets
the floor to leading order at an exponential build cost
(\Cref{thm:bhf-optimality}), and the perfect hash filter replaces the
Bloom filter's multiplicative $\log_2 e \approx 1.44$ overhead with an
additive addressing cost of at least $\log_2 e$ bits per element
(\Cref{prop:phf-space}), which wins at tight targets. On the FPR-zero
arm the BPHF, in the witness model, shifts the slot width from $\lceil
-\log_2 \fprate \rceil$ to $\lceil \log_2 n \rceil$ and the controlled
error from $\fprate$ to $\fnrate$. The two are the endpoints $(\fprate, 0)$ and
$(0, \fnrate)$ of the line; a structure that tolerates both errors can
interpolate between them, trading one axis against the other, and its
space floor with a bounded fraction of false negatives is the
lossy-dictionary bound of\index{space lower bound!with false negatives} \textcite[Prop.~1]{pagh-rodler-2001}. The space-accuracy duality of the chapter~6 geometry becomes
a deployment decision once the operating points are in hand:
false-negative-intolerant applications take an FNR-zero filter,
false-positive-intolerant ones the BPHF when the witness model holds.

\subsection*{Part V preview}

Part~V picks up two threads. The first is diagnostic: chapter~16
catalogues the structural ways a system fails to be Bernoulli, varying
rates, correlated errors, and adversarial inputs that break the
fixed-rate axioms, giving a checklist for when the framework applies
cleanly and when it needs modification. The witness model's dependence on
an unforgeable rank (\Cref{rem:witness-is-query-model}) is one entry in
that catalogue and the hinge to the second thread, which is generative:
chapters~17 and~18 develop \emph{cipher maps}\index{cipher map}, recasting the
perfect-hash-and-fingerprint architecture as the building block for
encrypted search, with the perfect hash replaced by a trapdoor whose
forward direction is public and reverse private. The FPR and FNR
guarantees of the present part carry into the cryptographic setting on
the same architectural template. \Cref{sec:bib-notes-ch15} attaches the
construction to its antecedents.
